\section{Constructor and Skolem sharing in the equational kernel}
\label{sec:e679}

The complete $688$-bit equational packet in the pinned 236-state report
has three closure-witness operations that can be identified with its
three syntax constructors.  The identifications are
\[
 \mathsf{Target}(b)=Q(b),\qquad
 \mathsf{Next}(b,z)=P(b,z),\qquad
 \mathsf{Mem}(b,z)=A(b,z).
 \tag{Fusion}
\]
They leave every equation tree and every variable occurrence unchanged.
Only nine operation occurrences in the final closure equation change
their names; three operation declarations disappear.  This gives a
complete $679$-bit packet $\mathcal E_{679}$ in the unchanged codec.

This is a specialization of the old signature, so merely renaming symbols
would not prove equiconsistency.  The substantive point is the upper
model below: a syntax code can have a prescribed membership type as well
as a prescribed syntactic role.  Infinitely many copies of each object
make the two requirements compatible.  The construction depends on the
fact that the interpreted matrices are equality-free.

\subsection*{The unchanged axioms and the lower implication}

Retain the first six equations of $\mathcal E_{688}$ literally: Sh1,
the native head and tail identities, NAND classes, atomic classes, and
the complemented binder.  In particular,
\begin{align*}
 V(P(x,s),d)&=x,&
 V(P(x,s),A(j,d))&=V(s,j),\\
 p_\varepsilon(s,P(a,b))
   &=N(p_\varepsilon(s,a),p_\varepsilon(s,b)),&
 p_\varepsilon(s,A(i,j))&=p_E(V(s,i),V(s,j)),\\
 p_\varepsilon(s,Q(a))
   &=N\bigl(p_\varepsilon(P(W(s,a),s),a),
             p_\varepsilon(P(z,s),a)\bigr).
\end{align*}
The complete closure equation is the old equation after (Fusion).
For its explanation, write
\[
 h=Q(b),\quad n_z=P(b,z),\quad m_z=A(b,z),\quad
 R_b(x,y)=\varepsilon(P(y,P(x,b)),V(b,d)).
\]
After taking Boolean truth values, the unchanged equation says
\[
 \mathsf{Bad}(b,z)\ \lor\ \mathsf{Hull}(b,x,y)
 \qquad\text{for all }b,z,x,y,
\]
where
\begin{align*}
 \mathsf{Bad}(b,z)
   &=R_b(h,n_z)\land E(m_z,n_z)\land\neg E(m_z,z),\\
 \mathsf{Hull}(b,x,y)
   &=\bigl(E(W(x,d),x)\lor\neg R_b(x,y)\lor E(y,h)\bigr)\\
   &\quad\land\bigl(\neg E(x,h)\lor\neg R_b(x,y)\lor E(y,h)\bigr)\\
   &\quad\land\bigl(\neg E(y,h)\lor R_b(P(b,y),y)\bigr).
\end{align*}
These displayed abbreviations are expanded in the supplied AST.

Every separating $\mathcal E_{679}$ algebra expands to an
$\mathcal E_{688}$ algebra by defining the three old operation symbols
according to (Fusion).  This is a literal substitution in each axiom;
the ground endpoints remain $N(d,d)$ and $d$.  The inherited
$\mathcal E_{688}$/raw-EPSC/ZFC calibration therefore proves
\[
 \mathcal E_{679}\text{ has a separating model}
       \ \Longrightarrow\ \operatorname{Con}(\mathrm{ZFC}).
\]
The derivation of $d$ as Boolean top from the equation packet also
survives by this same term substitution.  The reverse implication needs
the following construction.

\subsection*{A carrier with infinitely many copies}

Work in ZFC and assume $\operatorname{Con}(\mathrm{ZFC})$.  By first-order
completeness and downward L\"owenheim--Skolem, take a countably infinite
set-coded model $M$ of ZFC.  It may be nonstandard and non-well-founded
externally.  Fix a surjection
\[
             \tau:\mathbb N\longrightarrow |M|
\]
whose every fiber is infinite, and pull membership back to $\mathbb N$:
\[
          E(i,j)\quad\Longleftrightarrow\quad
                        M\models\tau(i)\in\tau(j).
\]
Call $\tau(i)$ the membership type of $i$.  Every equality-free
pure-membership formula is invariant under replacing any assignment by
one with the same types.  Structural induction proves this, with
surjectivity of $\tau$ treating the quantifiers.  Consequently all raw
EPSC instances hold on this carrier.  It need not satisfy Extensionality
with its literal equality, and no such assertion is required.

Set $d=0$.  We construct $P,A,Q$ with pairwise disjoint ranges, jointly
fresh outputs, and
\[
       P(i,j)>\max(i,j),\qquad A(i,j)>\max(i,j),\qquad Q(i)>i.
\]
They are injective in their complete argument lists.  Their output types
will be prescribed by the closure witnesses.  An infinite fiber always
has an unused element above any finite bound, so these requirements
leave a choice at every finite stage.

Interpret a $P$-code as a stream cell with its two unique children as
head and tail.  Every other element has head and tail zero.  The selector
index of $i$ is the number of successive outer layers $A(j,d)$; a
different outer shape has index zero.  Let $V(s,i)$ be the corresponding
stream coordinate.  The two native equations hold on all arguments,
including noncanonical streams and selectors.

Assign class $d$ and all nonconstructor numbers the constantly false
meaning.  Interpret a $P$-class as NAND of its two child classes, an
$A$-class as the indicated atomic membership formula, and a $Q$-class as
existential negation after prepending a stream coordinate.  Increasing
codes make this recursion well founded.  Each fixed class denotes a
finite equality-free pure-membership formula in finitely many stream
coordinates.  In particular, each fixed descriptor $b$ determines a
binary relation $R_b$ depending only on the types of its two arguments
and finitely many fixed parameter types.

\subsection*{Why the construction is not circular}

At stage $n$, perform the following finite batch in this order:
\begin{enumerate}
 \item Choose the value $Q(n)$.
 \item Choose every $P(b,z)$ with $\max(b,z)=n$.
 \item Choose every $A(b,z)$ with $\max(b,z)=n$.
\end{enumerate}
At each choice require the output to exceed $n$, its arguments, and all
outputs chosen previously.  Thus a number at most $n$ that is not already
a constructor output can never become one at a later stage.  Every
constructor decomposition at an index at most $n$ is already settled.

If $n$ is a $P$-code, its head $V(n,d)$ is strictly smaller than $n$;
otherwise that head is $d$.  Its fixed tail coordinates are likewise
smaller than $n$ or zero.  Therefore the finite pure-membership formula
represented by $R_n$ is determined before $Q(n)$ is chosen.  Although
its evaluation uses streams $P(y,P(x,n))$ for arbitrary $x,y$, decoding
such a stream always gives the literal coordinates $y,x,n$, independently
of the numerical labels or membership types of its new constructor
outputs.  Those future output choices cannot change the represented
formula.  This is the relevant distinction between syntax-code identities
and the types of the encoded objects.

For each descriptor $b$, mark one of two cases when $Q(b)$ is chosen.
If $R_b$ is locally subset bounded, choose an EPSC hull type $h_b$ and
give $Q(b)$ that type.  If it is not locally bounded, choose an input
type $h_b$ with unbounded output subsets and give $Q(b)$ that type.
Such choices exist in the pulled-back model, or equivalently in $M$.
This uses the inherited upper direction for raw EPSC only at a fixed
ordinary formula with a finite tuple of parameters.

When $P(b,z)$ is chosen, $Q(b)$ has already been chosen.  In the bounded
case, if $E(z,Q(b))$, prescribe its type to be an input witnessing
$R_b(P(b,z),z)$; if $z$ is outside the hull, prescribe any type.  EPSC
support supplies the witness.  In the unbounded case, prescribe an
output type satisfying
\[
  R_b(Q(b),P(b,z)),\qquad P(b,z)\not\subseteq_E z.
\]
The selected input type is unbounded, so such an output exists for every
$z$.  All copies of a prescribed type have the same required relation
values, permitting the fresh increasing code choice.

Finally, when $A(b,z)$ is chosen, $P(b,z)$ is already available.  In the
unbounded case choose its type so that
\[
                  E(A(b,z),P(b,z))\land\neg E(A(b,z),z).
\]
In the bounded case prescribe any type, because this witness is unused.
The order of the three parts of each batch supplies all dependencies.
Only finitely many codes have been allocated before every choice, so
none of the infinite type fibers can be exhausted.  This constructs total
operations meeting both their syntactic and their witness requirements.

\subsection*{Binder witnesses and Boolean values}

For $a\ne d$, define $W(s,a)$ as the least external-index witness
falsifying $\varepsilon(P(z,s),a)$ when such a witness exists, and zero
otherwise.  Class incidence is defined before $W$, so this introduces no
recursive dependence.  At $a=d$ the class is constantly false and the
complemented binder places no restriction on the witness value.  Retain
the old allocation
\[
 W(x,d)=\begin{cases}
    \min\{u:E(u,x)\},&\text{if $x$ has a member},\\
    0,&\text{otherwise}.
 \end{cases}
\]
All native, NAND, atomic, and binder identities hold.  The choices above
make the bad branch true for every $z$ in the unbounded case.  In the
bounded case, EPSC closure and support make every hull clause true;
the member selector makes its first clause vacuous at every nonempty
input and lets EPSC handle empty inputs.

Transport a countably infinite Boolean algebra to the carrier
$\mathbb N$ with top equal to $d$, interpret $N$ as NAND, and interpret
$p_E,p_\varepsilon$ by its top and bottom according to the two relations
just constructed.  Leave all native operations fixed.  The equation
packet holds and $N(d,d)\ne d$.  The Boolean algebra and membership types
are independent structures on the same carrier; no quotient of native
objects by a Boolean homomorphism is taken in this upper construction.

\paragraph{Theorem.}
With the inherited E688 calibration as an explicit dependency, ZFC proves
\[
 \mathcal E_{679}\not\vdash N(d,d)=d
       \quad\Longleftrightarrow\quad
                  \operatorname{Con}(\mathrm{ZFC}).
\]
Indeed, the preceding construction gives a separating algebra from
consistency; the literal expansion gives the reverse implication.
The usual quotient of ground terms by the generated equational
congruence connects separating models to nonderivability.  This is an
ordinary proof formalizable in ZFC, not a newly checked Lean/Rocq theorem
or an internal PA proof-code transformer.

\subsection*{Exact cost and limitations}

The new signature is
\[
       d^0,N^2,V^2,P^2,A^2,Q^1,W^2,p_E^2,p_\varepsilon^2.
\]
There are still seven equations, four local variable slots, 153 equation
tree nodes, and four endpoint nodes.  Nine symbols plus four variables
still require a four-bit opcode.  The complete charged count is
\[
 \underbrace{8+3\cdot9+8}_{43\text{ header}}
 +\underbrace{7+1+4\cdot153}_{620\text{ equations}}
 +\underbrace{4\cdot4}_{16\text{ endpoints}}
 =679.
\]
The nine-bit saving consists exactly of the three removed arity entries;
it does not shorten the equation payload under this fixed-width codec.
This is a small numerical gain and a more general structural sharing
principle.  It supplies no Turing-machine state saving on its own.
An actual binary implementation must still charge a ground-rewrite
checker, fair enumerator, complete packet writer, and blank startup.

The accompanying builder pins and reconstructs E688, checks every
corresponding AST node, confirms the unchanged first six equations and
endpoints, round-trips E679, and rejects all strict packet truncations.
Its checks establish literal transcription and encoding facts; the
infinite-model argument is the mathematical justification for the
specialization.  Equality-bearing matrix semantics would invalidate the
type-invariance step, so the fusion must not be imported unchanged into
the original equality-bearing C/G machine calculus.

\paragraph{An optional eight-operation specialization.}
One can also eliminate the remaining binder-witness symbol by
\[
                    W(s,a)=A(Q(a),s).
\]
This further specialization has a compatible upper model.  A descriptor
$b=Q(a)$ is not a $P$-code, so $V(b,d)=d$ and $R_b$ is constantly false.
It is therefore in the bounded case with an empty hull, and all witness
slots $A(b,z)$ are free of closure-witness constraints.  Use precisely
those free slots for the binder counterexamples, retaining the member
selector property when $a=d$: a fresh copy of the type of a selected
member still belongs to the nonempty input class. Empty inputs impose no
membership-witness requirement here.  Injectivity of $Q$ makes the allocation
unambiguous.  At the stage when $A(Q(a),s)$ is chosen, $a<Q(a)$, so the
class $a$ and the finite coordinates of stream $s$ are already fixed;
type invariance permits any sufficiently fresh copy of the selected
witness type.  Thus this use of $A$ introduces no circular dependence
between binder choices and class meanings.

There are two $W$ occurrences.  Replacing each by $A(Q(a),s)$ adds one
AST node while removing one three-bit arity declaration.  The same
fixed-width codec would therefore charge $679+2\cdot4-3=684$ bits.
It reduces the signature from nine to eight operations but increases
the complete packet under this codec.  The supplied primary packet
remains E679; the eight-operation variant is a proved alternative
specialization, not an additional claimed state reduction.
